\centerline{\bf \Huge Клеточки))}

\begin{flushright}
        \scriptsize{}
\end{flushright}

\problem{1} В клетках таблицы $5 \times 5$ расставлены натуральные числа так, что все десять сумм этих чисел в строках и столбцах таблицы различны. Найдите наименьшее возможное значение суммы чисел во всей таблице.

\problem{2} На некоторых клетках доски $10 \times 10$ стоят шашки. Клетка называется красивой, если на горизонтали, проходящей через эту клетку, стоит нечетное число шашек, и на вертикали, проходящей через ту же клетку, тоже стоит нечетное число шашек. Может ли на доске оказаться ровно 42 красивые клетки?

\problem{3} В клетки квадрата $10 \times 10$ расставили числа $1,2,3, \ldots, 100$ (в каждую клетку одно число). Оказалось, что нет двух соседних по стороне чисел, разность которых делится на 3. Затем посчитали количество трёхклеточных уголков внутри квадрата, сумма чисел в клетках которых делится на 3 . Какое наибольшее число могло получиться?

\problem{4} В каждой клетке таблицы $100 \times 100$ записано натуральное число. В каждой строке имеется по крайней мере 10 различных чисел, а в каждых четырех последовательных строках не более 15 различных чисел. Какое наибольшее количество различных чисел может быть в таблице?

\problem{5} В каждой клетке таблицы $N \times N$ записано число. Назовём клетку хорошей, если сумма чисел строки, содержащей эту клетку, не меньше, чем сумма чисел столбца, содержащего эту клетку. Найдите наименьшее возможное количество хороших клеток.

\problem{6} Клетки таблицы $200 \times 200$ окрашены в черный и белый цвета так, что черных клеток на 404 больше, чем белых. Докажите, что найдется квадрат $2 \times 2$, в котором число белых клеток нечетно.
